\documentclass{article}
\usepackage{amsmath,amssymb,amsthm,hyperref}
\newtheorem{theorem}{Theorem}
\title{Exact limit-preservation boundary for native finite-witness wiring}
\date{16 November 2026}
\begin{document}\maketitle
\paragraph{Audit and scope.}
This addendum concerns exactly the counting-witness action of
\texttt{Finite\_Witness\_Wiring\_Module.tex}, not every semantics of spans or
wiring. Assume the context and writer coefficient categories are symmetric
monoidal and finitely complete. Use all small families in each fixed-interface
fibre, retaining raw zero-weight witnesses and entire global writer grades.
Finite-set interfaces and finite wiring apices remain mandatory. The statement
classifies when the existing canonical comparison is invertible; it does not
supply a new universal completion or a historical novelty claim.

For $w=(X\xleftarrow{s}U\xrightarrow{t}Y)$ write
$A_w:\operatorname{Fam}(\mathsf S_X)\to\operatorname{Fam}(\mathsf S_Y)$
for the componentwise action. A native system $(W,p,a)$ is sent to the raw
carrier $W\times_XU$, with weight $p(b)$ at $(b,z)$ and output $t(z)$.
Coefficients and its entire global writer grade are unchanged. For each finite
diagram $E$, let
\[
 \kappa_{w,E}:A_w(\lim E)\longrightarrow\lim A_w(E)
\]
be the map determined by the acted-on cone. Limits here are shared-boundary
fibre limits, not arbitrary marginal products at unrelated interfaces.

\begin{theorem}[Sharp classification]
For a finite span $w$ the following statements hold.
\begin{enumerate}
\item The terminal-object comparison is invertible if and only if both $s$
and $t$ are bijections.
\item Every binary-product comparison is invertible if and only if $s$ is a
bijection and $t$ is injective. These conditions are also equivalent to
invertibility of $\kappa_{w,E}$ for every nonempty finite diagram $E$.
\item Every finite-limit comparison is invertible if and only if both $s$ and
$t$ are bijections. In that case $A_w$ is an equivalence given by relabelling
along the bijection $t s^{-1}$, up to its specified raw tuple isomorphism.
\end{enumerate}
Here a nonempty diagram has at least one object; it may be disconnected and
its limit family may be empty.
\end{theorem}

\begin{proof}[Terminal comparison]
The terminal family at $X$ has one component for every finite nonnegative
measure $\mu$ on $X$. Each component has raw carrier $X$, boundary identity,
terminal context coefficients and terminal global writer grade. The acted-on
component has raw carrier $U$, law $\mu(s(z))$, boundary $t$, and the same
terminal data. Its map to the terminal family at $Y$ has raw function $t$ and
family-index function
\[
 \Phi(\mu)=t_*(\mu\circ s).
\]
An isomorphism of families must have a bijective index function and a native
isomorphism in each component. In particular its raw maps must be bijective.
There is always at least the zero measure component, including when $X$ is
empty. Thus invertibility forces $t$ bijective. Under this hypothesis,
$\Phi$ is bijective exactly when restriction
$\mathbb R_{\ge0}^X\to\mathbb R_{\ge0}^U$ along $s$ is bijective. If $s$
misses a point, measures differing only there have the same restriction. If
$s$ identifies two points of $U$, a function with unequal values there cannot
be a restriction. Conversely a bijective $s$ gives a bijective restriction.
All coefficient and grade components are then the canonical terminal
isomorphisms, establishing sufficiency as well as necessity.
\end{proof}

\begin{proof}[Necessity for binary products]
Suppose all binary-product comparisons are invertible. First, $s$ must be
surjective. If $x\in X$ is omitted, take two singleton systems at $x$, of
masses one and two, with unit coefficients and writer grades. Their
shared-boundary product family is empty, since a joint measure cannot have
both total masses. Both acted-on systems have empty raw carrier and zero
measure. Their product family is nonempty: the empty raw carrier with zero
measure and the full product writer grade is a component. The comparison
from an empty family cannot be invertible.

Next suppose $F=s^{-1}(x)$ contains at least two points. If $t$ identifies
two points of $F$, use two singleton mass-one systems at $x$. Their original
product has one law and one raw witness; after action its raw carrier is $F$.
The target product raw carrier is
\[
 \{(z,z')\in F^2:t(z)=t(z')\}.
\]
The comparison is the diagonal map and selects the diagonal counting law.
Its raw map misses an off-diagonal point. That point remains a genuine raw
witness even though its weight in the selected law is zero. Thus the selected
component map itself is not an isomorphism.

If instead $t$ is injective on $F$, use two fair-bit systems at $x$. A source
product component is a bit coupling $\lambda$, and its image uses that same
coupling separately at every $z\in F$. The target product has raw carrier
$\{0,1\}^2\times F$ and admits a distinct fair-bit coupling at every $z$:
take the diagonal one at one point and the antidiagonal one at another, with
either choice at remaining points. Its marginals are exactly the two acted-on
fair-bit laws. It is not in the image of the family-index function. Hence the
comparison again fails to be an isomorphism. These cases prove injectivity of
$s$.

With $s$ now bijective, suppose $t$ identifies distinct $z,z'$. Let
$x=s(z)$ and $x'=s(z')$. Singleton mass-one states at these distinct boundary
values have empty shared-boundary product. Their images are singleton
mass-one states at the same output boundary and have a nonempty product.
This again contradicts invertibility. Consequently $t$ must be injective.
None of these tests uses a coefficient shadow: the chosen units and their
actual finite products supply the required coefficient and full-grade data.
\end{proof}

\begin{proof}[Sufficiency and remaining implications]
If $s$ is bijective, projection $W\times_XU\to W$ is a native isomorphism
identifying $A_w$ with relabelling of the boundary along $f=ts^{-1}$. If also
$t$ is injective, so is $f$. Consider any nonempty finite diagram of families.
The coherent tuples of family indices are unchanged by relabelling. For each
such tuple, equality of output boundary values $f(a_i(b_i))$ is equivalent
to equality of the original values $a_i(b_i)$. Thus its complete compatible
raw tuple carrier is identical before and after relabelling. The diagrams of
context coefficients and entire global writer grades are identical, and the
admitted finite joint laws with their prescribed marginals are identical.
The canonical comparison is therefore an isomorphism on family indices and
on all native data. This includes inconsistent diagrams with no admissible
indices or laws. Nonemptiness matters: in the empty diagram the terminal raw
carrier changes from $X$ to $Y$, and arbitrary laws outside $f(X)$ have no
preimage.

Preservation of all nonempty finite limits implies binary-product preservation
by taking a discrete two-object diagram; the preceding necessity argument
therefore establishes the stated equivalence. Combining the nonempty case
with the terminal classification proves the finite-limit classification.
If $s,t$ are bijective, inverse relabelling gives the stated equivalence.
\end{proof}

\paragraph{Empty-interface edge.}
If $X=\varnothing$, then $U=\varnothing$. The empty $s$ is bijective and the
empty $t$ is injective for every $Y$. The action therefore preserves every
nonempty finite limit, including the full unused global writer-grade data
carried by empty raw systems. It preserves a terminal object only if
$Y=\varnothing$. This edge is included by the theorem rather than excluded
by an implicit positive-mass assumption.

\paragraph{Implication for feedback.}
The equality-feedback span from the main manuscript generally has a
non-surjective left leg (nonclosed boundary assignments are excluded), and
its right leg can identify distinct closed loop values. The theorem explains
why strong joint-limit preservation is exceptional, while the canonical
directed comparison remains valid. It does not identify this finite
instantaneous closure with recurrent or normalized feedback semantics.
\end{document}
